% docs/slides/preamble.tex -- 五份 deck 共享
\documentclass[aspectratio=169,11pt]{beamer}

% ==== 中文与字体 ====
\usepackage[UTF8,fontset=none]{ctex}
\usepackage{fontspec}

\IfFontExistsTF{PingFang SC}{%
  \setCJKmainfont{PingFang SC}[BoldFont=PingFang SC,ItalicFont=PingFang SC,AutoFakeSlant=0.15]
  \setCJKsansfont{PingFang SC}[ItalicFont=PingFang SC,AutoFakeSlant=0.15]
  \setCJKmonofont{PingFang SC}
}{%
  \IfFontExistsTF{Source Han Sans SC}{%
    \setCJKmainfont{Source Han Sans SC}
    \setCJKsansfont{Source Han Sans SC}
  }{%
    \IfFontExistsTF{Noto Sans CJK SC}{%
      \setCJKmainfont{Noto Sans CJK SC}
      \setCJKsansfont{Noto Sans CJK SC}
    }{%
      \setCJKmainfont{Heiti SC}
      \setCJKsansfont{Heiti SC}
    }
  }
}

\IfFontExistsTF{Helvetica Neue}{\setmainfont{Helvetica Neue}}{%
  \IfFontExistsTF{Inter}{\setmainfont{Inter}}{}}
\IfFontExistsTF{JetBrains Mono}{\setmonofont{JetBrains Mono}[Scale=0.85]}%
  {\setmonofont{Menlo}[Scale=0.85]}

% ==== 颜色 ====
\usepackage{xcolor}
\definecolor{cMain}{HTML}{1F4E79}
\definecolor{cAccent}{HTML}{E07A1F}
\definecolor{cGray}{HTML}{595959}
\definecolor{cMute}{HTML}{F4F4F4}
\definecolor{cOK}{HTML}{2E7D32}
\definecolor{cBad}{HTML}{B23A3A}
\definecolor{cWarn}{HTML}{C8A300}

% ==== 主题 ====
\usepackage{tcolorbox}
\usetheme{metropolis}
\metroset{progressbar=frametitle,numbering=fraction,block=fill}
\setbeamercolor{frametitle}{bg=cMain,fg=white}
\setbeamercolor{progress bar}{fg=cAccent,bg=cMute}
\setbeamercolor{alerted text}{fg=cAccent}
\setbeamercolor{title}{fg=cMain}

% metropolis 默认尝试 Fira Sans，未安装时 fallback 到 Latin Modern Sans。
% 这里在主题加载后显式覆盖，强制使用 macOS 上一定存在的 Helvetica Neue，
% 跟中文 PingFang 风格一致；找不到则按顺序尝试 Inter / Arial。
\IfFontExistsTF{Helvetica Neue}{%
  \setsansfont{Helvetica Neue}%
}{%
  \IfFontExistsTF{Inter}{\setsansfont{Inter}}{%
    \IfFontExistsTF{Arial}{\setsansfont{Arial}}{}}%
}

% ==== TikZ ====
\usepackage{tikz}
\usetikzlibrary{positioning,arrows.meta,fit,shapes.geometric,calc,%
  decorations.pathmorphing,backgrounds}

\tikzset{
  node-box/.style={draw=cMain,thick,rounded corners=2pt,
    minimum height=8mm,minimum width=18mm,align=center,font=\small},
  node-state/.style={draw=cMain,thick,circle,minimum size=14mm,
    align=center,font=\small},
  node-ok/.style={draw=cOK,thick,fill=cOK!10},
  node-warn/.style={draw=cWarn,thick,fill=cWarn!15},
  node-bad/.style={draw=cBad,thick,fill=cBad!10},
  % 色盲友好的状态机三态：蓝 (正常/Closed) / 橙 (告警/Open) / 灰 (中间/HalfOpen)
  % 避免 red-green 二元区分，改用 hue + lightness 双通道。
  node-state-normal/.style={draw=cMain,thick,fill=cMain!12},
  node-state-alert/.style={draw=cAccent,thick,fill=cAccent!18},
  node-state-transition/.style={draw=cGray,thick,fill=cGray!18},
  arrow/.style={-{Stealth[length=2mm]},thick,draw=cMain},
  arrow-async/.style={-{Stealth[length=2mm]},thick,dashed,draw=cGray},
  arrow-bad/.style={-{Stealth[length=2mm]},thick,draw=cBad},
  arrow-alert/.style={-{Stealth[length=2mm]},thick,draw=cAccent},
  lifeline/.style={dashed,thick,draw=cGray},
}

% ==== 代码 ====
\usepackage{listings}
\lstdefinelanguage{Go}{
  morekeywords={break,case,chan,const,continue,default,defer,else,
    fallthrough,for,func,go,goto,if,import,interface,map,package,range,
    return,select,struct,switch,type,var,nil,true,false,iota},
  morekeywords=[2]{string,int,int64,uint,uint64,bool,byte,error,context},
  sensitive=true,
  morecomment=[l]//,morecomment=[s]{/*}{*/},
  morestring=[b]",morestring=[b]`,
}
\lstdefinelanguage{LuaRedis}{
  morekeywords={and,break,do,else,elseif,end,false,for,function,if,in,
    local,nil,not,or,repeat,return,then,true,until,while},
  morekeywords=[2]{redis,call,KEYS,ARGV,tonumber,HGET,HSET,GET,SET,
    DECR,DECRBY,INCR,INCRBY,EXPIRE,PEXPIRE,ZADD,ZCARD,ZREMRANGEBYSCORE,
    EVAL,EVALSHA},
  sensitive=true,
  morecomment=[l]{--},
  morestring=[b]",morestring=[b]',
}

\lstdefinestyle{base}{
  basicstyle=\ttfamily\scriptsize,
  numbers=left,numberstyle=\tiny\color{cGray},
  keywordstyle=\color{cMain}\bfseries,
  keywordstyle=[2]\color{cAccent},
  stringstyle=\color{cOK},
  commentstyle=\color{cGray}\itshape,
  backgroundcolor=\color{cMute},
  frame=single,framerule=0pt,
  showstringspaces=false,
  breaklines=true,
  tabsize=2,
  columns=fullflexible,
}
\lstset{style=base}

% ==== 三线表与附录页码 ====
\usepackage{booktabs}
\usepackage{appendixnumberbeamer}

% ==== 页脚来源标注 ====
\newcommand{\srcnote}[1]{{\color{cGray}\scriptsize\itshape #1}}

% 关键论点框
\newcommand{\keypoint}[1]{%
  \begin{tcolorbox}[colback=cMute,colframe=cMain,boxrule=0.5pt,arc=2pt]%
  #1\end{tcolorbox}}


\title{商品展示的业务依据}
\subtitle{转化漏斗最宽一跳、首屏 1 秒、三层缓存与深分页反模式}
\author{gomall · 业务侧 deck}
\date{}

\begin{document}

\maketitle

\begin{frame}{今日大纲}
\tableofcontents
\end{frame}

% ============================================================
\section{商品展示在业务全景中的位置}
% ============================================================

\begin{frame}{商品展示是什么角色的痛点}
\begin{tabular}{@{}p{16mm}p{42mm}p{42mm}@{}}
\toprule
角色 & 关心什么 & 商品展示影响 \\
\midrule
C 端用户 & 搜得到 / 看得清 / 不卡 & 详情页 RT 100ms = 跳出率 +2\%；首屏 \textgreater 3s 流失 90\% \\
商家     & 我上架的 SKU 多久能在首页可见 & 上架到可见 SLA：详情 60s / 列表 30s / 类目 5min \\
运营     & 轮播/类目导流是不是有效 & 轮播加载慢 1s = 首页 GMV $-$2\%；命中率掉档直接打 DAU \\
客服     & 用户投诉"显示老价格"怎么解释 & 双删 500ms / HTTP 60s / outbox 落单分别对应三套话术 \\
SRE      & 双 11 0 点能不能扛 100k QPS & 99\% Redis 命中 = DB 实际只看 320 QPS，单机够用 \\
\bottomrule
\end{tabular}
\vspace{2mm}
\small 商品展示是\textbf{转化漏斗最宽的一跳}：用户必然经过、必然以毫秒计耐心、必然有商家变更冲突。这一节用业务侧 70\% 的篇幅讲"为什么这么做"，代码 30\%。
\srcnote{README.md \S 业务全景 / 用户旅程 deck 02 节点}
\end{frame}

\begin{frame}{用户旅程视图：详情页是漏斗的腰}
\begin{tikzpicture}[node distance=2mm and 4mm,font=\tiny]
  \node[node-box,minimum width=18mm,node-ok] (login) {登录};
  \node[node-box,minimum width=22mm,node-ok,right=of login] (browse) {\textbf{浏览/详情}};
  \node[node-box,minimum width=18mm,node-ok,right=of browse] (cart) {加购};
  \node[node-box,minimum width=18mm,node-ok,right=of cart] (order) {下单};
  \node[node-box,minimum width=18mm,node-ok,right=of order] (pay) {支付};
  \node[node-box,minimum width=18mm,node-ok,right=of pay] (ship) {履约};
  \draw[arrow] (login) -- (browse);
  \draw[arrow] (browse) -- (cart);
  \draw[arrow] (cart) -- (order);
  \draw[arrow] (order) -- (pay);
  \draw[arrow] (pay) -- (ship);
  \node[below=8mm of login,font=\tiny] {DAU 100\%};
  \node[below=8mm of browse,font=\tiny\bfseries\color{cAccent}] {\textasciitilde 80\%};
  \node[below=8mm of cart,font=\tiny] {\textasciitilde 15\%};
  \node[below=8mm of order,font=\tiny] {\textasciitilde 5\%};
  \node[below=8mm of pay,font=\tiny] {\textasciitilde 4.5\%};
  \node[below=8mm of ship,font=\tiny] {\textasciitilde 4.5\%};
\end{tikzpicture}
\vspace{3mm}
\begin{itemize}
  \item 详情页一次抖动 = 漏斗腰部塌陷 = 下游所有环节 GMV 同比降。
  \item gomall 把详情 / 列表 / 类目 / 轮播 4 条最宽接口都挂了 HTTP cache + Redis cache。
  \item 同样的优化用在"加购"以下没意义：那里量已经降到 15\%，瓶颈是库存与事务一致性而不是缓存命中率。
\end{itemize}
\srcnote{routes/router.go:39--45 四个公开 GET 接口集中挂 HTTPCache}
\end{frame}

\begin{frame}{业务侧 SLO：详情页背的指标}
\footnotesize
\begin{tabular}{@{}lllr@{}}
\toprule
指标 & 业务等级 & 目标 & 实测 \\
\midrule
详情页 p99 延迟        & P1 核心读 & \textless 200ms & 3ms (p95) \\
列表/类目 p99 延迟     & P3 信息类 & \textless 1s    & 列表 2.5s (反例) \\
可用率                 & P1        & 99.9\%          & N/A \\
Redis 命中率           & 内部 KPI  & $\ge$ 99\%      & N/A \\
端到端 DB 流量占比     & 内部 KPI  & $\le$ 0.5\%     & 公式推 0.32\% \\
商家上架 $\to$ 可见 SLA & 业务对齐  & 60s             & = HTTP max-age \\
商家改价 $\to$ 可见 SLA & 业务对齐  & 1s              & 双删 500ms + TTL \\
\bottomrule
\end{tabular}
\vspace{2mm}

SLO 是\textbf{业务和技术之间的契约}：超了赔运营预算，没超就接受延迟。商品列表 24.5 RPS / p95 2.5s \textbf{不算事故}（P3 上限是 1s 但当前 QPS 远低于阈值），但已经进入"应该排期修"的状态。
\srcnote{README.md \S 业务承诺(SLO) P1/P3 行 + stressTest/REPORT.md 链路压测表}
\end{frame}

\begin{frame}{业务边界：这一节不解决什么}
\begin{itemize}
  \item \textbf{不做 SKU 多规格详情}：当前商品 model 是单 SKU，没有"颜色/尺码"维度。
  \item \textbf{不做视频 / 直播流}：详情页只有 \texttt{ImgPath} 单图，视频解码 / 切片 / CDN 调度独立链路。
  \item \textbf{没有 AB 测试框架}：双 max-age 切流量、内容个性化、灰度发布都需要外置 AB 平台。
  \item \textbf{没有个性化推荐}：详情页底部"猜你喜欢"位是空的，路线图依赖召回 + 排序服务。
  \item \textbf{商品搜索召回}：见 deck 03（ES + Milvus hybrid），不在本 deck。
  \item \textbf{评价 / 问答 / 售后入口}：见 deck 09，路线图阶段。
\end{itemize}
\vspace{1mm}
\small 摊牌不丢人。把不做什么写清楚比把做了什么写得花哨更值钱：让面试官 / 商家 / 运营知道边界在哪里。
\srcnote{README.md \S 业务边界(不做什么) + docs/architecture/feature-matrix.md}
\end{frame}

% ============================================================
\section{首屏 1 秒决定订单}
% ============================================================

\begin{frame}{这一节回答四个问题}
\begin{itemize}
  \item 商品详情页是首屏黄金 1 秒，缓存命中 vs 回源 DB 在业务上差多少。
  \item Cache Aside 读路径 + 双删写路径，为什么\textbf{业务必须接受}最终一致而不是强一致。
  \item HTTP cache (ETag + Cache-Control) 哪几个接口能挂、TTL 怎么定。
  \item \texttt{product/list} 跑 24.5 RPS / p95 2.5s 是怎么从业务侧解读为反例的。
\end{itemize}
\srcnote{对应代码 routes/router.go:39--45 / middleware/httpcache.go / repository/cache/product.go / internal/product/service.go}
\end{frame}

\begin{frame}{首屏延迟与转化率：业务侧的硬性约束}
\begin{tabular}{@{}lll@{}}
\toprule
首屏延迟 & 跳出率上升 & 业内常见处理 \\
\midrule
$\le$ 1s   & 基线        & 不动 \\
1 -- 3s    & +30\%       & 进入优化排期 \\
3 -- 5s    & +90\%       & 列入 P0 故障复盘 \\
$\ge$ 5s   & 用户基本走完  & 直接告警呼叫 \\
\bottomrule
\end{tabular}
\vspace{3mm}
\small Akamai 和 Google 在 2017 / 2018 公开的转化率曲线，是后端做缓存预算的业务依据。\\
gomall 的 \texttt{product/show} 压测 p95 = 3.01ms，留出 300 倍余量给业务侧的图片加载、首屏渲染。
\srcnote{stressTest/REPORT.md \S 各链路压测结果（行 product/show）}
\end{frame}

\begin{frame}{命中缓存 vs 回源 DB 的业务收益}
\begin{tabular}{@{}lrrr@{}}
\toprule
路径 & 单次耗时 & 956K 行表能承担 QPS & 业务侧含义 \\
\midrule
浏览器 304 & $\sim$ 30ms (RTT)  & 不打到源站 & 流量被 CDN/浏览器卸掉 \\
Redis 命中 & $\sim$ 1ms          & 数万级    & 单机 Redis 上限 \\
PK 回源 DB & 3.01ms (p95)        & 6 万级    & 与 ping 同量级 \\
全表扫     & 2.5s   (p95)        & 24.5      & 直接告警 \\
\bottomrule
\end{tabular}
\vspace{2mm}
\small "回源" 不可怕，"扫表" 才可怕。gomall \texttt{/product/show} 走主键 + 缓存兜底，业务侧从来不是瓶颈；
\texttt{/product/list} 全表 COUNT 才是要修的那个。
\srcnote{stressTest/REPORT.md \S 1 基线 vs 业务接口；\S 3 商品列表反模式}
\end{frame}

% ============================================================
\section{读路径：Cache Aside 是怎么挡住 DB 的}
% ============================================================

\begin{frame}{业务场景：商家上架 1 万 SKU 后 DB 顶不住的剧本}
\textbf{角色}：商家「数码玩家」周末批量上架 1 万个手机壳 SKU；运营把头部 200 个挂上首页推荐位。
\vspace{1mm}
\begin{itemize}
  \item \textbf{周一 10:00}：推荐位开始引流，C 端用户点击进详情页。每 SKU 详情页 \texttt{product:detail:\{id\}} 第一次访问 = cache miss = 回源 DB。
  \item \textbf{若没有 SETNX 单飞}：200 个热 SKU $\times$ 50 个并发用户 = 一万次并发 SELECT 同时打 DB；MySQL 连接池打满，整站详情页 503。
  \item \textbf{客服 10:05 开始接投诉}："首页推荐点不开"、"详情页一直转圈"；客服只能挂 70002 话术"服务暂忙"，但根因其实是缓存击穿。
  \item \textbf{有 SETNX 单飞}：同一商品最多 1 个请求打 DB，其余 sleep 50ms 后命中回填值；DB 看到的额外负载 = 200 QPS（每热 SKU 1 次），可忽略。
\end{itemize}
\keypoint{Cache Aside 的真正价值不是"快"，是"把商家行为爆点和 DB 解耦"。}
\srcnote{internal/product/service.go:42--68 / repository/cache/product.go:55--58 TryProductLock}
\end{frame}

\begin{frame}{Cache Aside 读路径}
\begin{tikzpicture}[node distance=8mm and 14mm]
  \node[node-box] (cli) {Client};
  \node[node-box,right=of cli,node-ok] (cache) {Redis\\product:detail:\{id\}};
  \node[node-box,right=of cache] (db) {MySQL\\product PK};
  \node[font=\tiny,below=2mm of cache] (lab) {未命中};
  \draw[arrow] (cli) -- node[above,font=\tiny]{GET /product/show} (cache);
  \draw[arrow] (cache) -- node[above,font=\tiny]{SELECT * WHERE id=?} (db);
  \draw[arrow] (db.south) -- ++(0,-6mm) -| node[pos=0.25,below,font=\tiny]{回填} (cache.south);
  \draw[arrow,bend left=30] (cache.north) to node[above,font=\tiny]{命中: JSON} (cli.north);
\end{tikzpicture}
\vspace{1mm}
\begin{itemize}
  \item 读：先读缓存；缺则回源，再写回缓存。
  \item gomall 实现里多一道 \texttt{SETNX} 回源锁：缓存击穿时 \textbf{只放一个请求} 打 DB，其它请求等 50ms 再读一次。
  \item 业务结果：商品详情压测 62K RPS，与 \texttt{/ping} 基线几乎重合。
\end{itemize}
\srcnote{internal/product/service.go:42--68 ProductShow 三段式}
\end{frame}

\begin{frame}[fragile]{读路径关键代码}
\begin{lstlisting}[language=Go,firstnumber=42,
  caption={\srcnote{internal/product/service.go:42--68 ProductShow}}]
cached := &ProductResp{}
if err := cache.GetProductDetail(ctx, req.ID, cached); err == nil {
    return cached, nil                          // L1: 命中, ~1ms
}
locked, _ := cache.TryProductLock(ctx, req.ID)  // L2: SETNX 抢回源权
if !locked {
    time.Sleep(50 * time.Millisecond)           // L2': 让锁持有者回填
    if err := cache.GetProductDetail(ctx, req.ID, cached); err == nil {
        return cached, nil
    }
} else { defer cache.UnlockProduct(ctx, req.ID) }
pResp, err := s.loadProductFromDB(ctx, req.ID)  // L3: 锁持有者打 DB
if err != nil { return nil, err }
_ = cache.SetProductDetail(ctx, req.ID, pResp)  // L4: 回填, TTL=10min
return pResp, nil
\end{lstlisting}
\end{frame}

\begin{frame}{读路径四步的业务含义}
\begin{description}
  \item[L1 命中] 99\% 的请求走这里。Redis $\sim$ 1ms，对业务侧等价于"没有数据库"。
  \item[L2 SETNX] 缓存失效/被双删的瞬间，防止"100 个并发同时打 DB"导致雪崩，业务侧叫\textbf{缓存击穿}。
  \item[L2'\,等] 没抢到锁的请求 sleep 50ms 再读一次。读到就走 L1；读不到再硬回源（兜底，不能因为锁丢请求）。
  \item[L4 回填] TTL = \texttt{ProductDetailTTL = 10min}，这是 Redis 层；HTTP 层是 60s。两层 TTL 不一样，下一节解释。
\end{description}
\srcnote{repository/cache/product.go:14 ProductDetailTTL = 10*time.Minute}
\end{frame}

% ============================================================
\section{写路径：双删的业务必要性}
% ============================================================

\begin{frame}{业务取舍：用户多久能看到改价}
\footnotesize
\textbf{场景}：商家在后台把"小米 14 Pro"从 ¥5499 改成 ¥4999，限时秒杀 1 小时。
\vspace{1mm}
\begin{center}
\begin{tabular}{@{}llp{5.8cm}@{}}
\toprule
方案 & 用户可见延迟 & 业务后果 \\
\midrule
强一致 (2PC / 同步写 cache) & $<$ 50ms       & 商家点"保存"按钮卡 200ms；DB / Redis 任一抖动 = 改价失败 \\
\textbf{最终一致 (双删)}    & $\le$ 500ms    & gomall 当前。商家无感，用户最多看到 0.5s 旧价 \\
TTL 自然过期 (无双删)       & $\le$ 10min    & 不可接受：秒杀已经开始 5 分钟，用户还看到老价 \\
\bottomrule
\end{tabular}
\end{center}
\vspace{2mm}

\textbf{为什么不接受 10 分钟旧价}：限时秒杀按分钟计价值；用户看到老价 = 漏发优惠 = 客诉 + 平台担责。\\
\textbf{为什么 500ms 可接受}：和产品 / 运营当面对齐："改价后 1 秒内可能仍显示老价，刷新即可"——写进客服话术。\\
\textbf{为什么不做强一致}：商家批量改 1000 个 SKU，每个卡 200ms = 后台转 3 分钟，完全不能用。
\srcnote{internal/product/service.go:267--292 ProductUpdate + repository/cache/product.go:64--74 DoubleDeleteAsync}
\end{frame}

\begin{frame}{为什么不能"写库 + 写缓存"}
\begin{itemize}
  \item 直觉做法：商家改价 $\to$ \texttt{UPDATE product} $\to$ \texttt{SET cache}。
  \item 业务现实里两个并发请求会：
  \begin{enumerate}
    \item T1 改价 100 $\to$ 99
    \item T2 改价 100 $\to$ 88
    \item T2 先到 DB（写 88），T1 后到 DB（覆盖为 99）
    \item T1 后写 cache (99)，T2 先写 cache (88)
    \item DB = 99，Cache = 88 —— \textbf{永久不一致}
  \end{enumerate}
  \item 客诉版本："为什么列表写 88，详情写 99？"
\end{itemize}
\keypoint{业务侧不能依赖"写缓存"，只能依赖"删缓存"。}
\srcnote{典型双写不一致问题；详见 Facebook TAO 论文 \S 4}
\end{frame}

\begin{frame}{延迟双删：第二刀挡哪个窗口}
\begin{tikzpicture}[node distance=10mm]
  \foreach \name/\x in {写者/0, Cache/40, 读者/80, DB/120} {
    \node[node-box,minimum width=14mm] at (\x mm,0) (\name) {\name};
    \draw[lifeline] (\x mm,-2mm) -- (\x mm,-60mm);
  }
  \draw[arrow] (0,-8mm)   -- node[above,font=\tiny]{Del cache} (40mm,-8mm);
  \draw[arrow] (80mm,-14mm)  -- node[above,font=\tiny]{Get cache (miss)} (40mm,-14mm);
  \draw[arrow] (80mm,-20mm)  -- node[above,font=\tiny]{回源 (拿到旧值)} (120mm,-20mm);
  \draw[arrow] (0,-26mm)   -- node[above,font=\tiny]{Update DB (新值)} (120mm,-26mm);
  \draw[arrow-bad] (80mm,-32mm) -- node[above,font=\tiny\color{cBad}]{回填 cache = 旧值} (40mm,-32mm);
  \draw[arrow-alert] (0,-44mm) -- node[above,font=\tiny\color{cAccent}]{500ms 后 Del cache (第二刀)} (40mm,-44mm);
  \node[font=\tiny,align=center] at (60mm,-54mm) {第二刀清掉读者写回的旧值\\下次读触发 L3 回源拿新值};
\end{tikzpicture}
\srcnote{repository/cache/product.go:64--74 DoubleDeleteAsync (interval=500ms)}
\end{frame}

\begin{frame}[fragile]{写路径关键代码：删-写-延迟删}
\begin{lstlisting}[language=Go,firstnumber=271,
  caption={\srcnote{internal/product/service.go:271--292 ProductUpdate}}]
product := &Product{ /* 价格/标题等字段 */ }
_ = cache.DelProductDetail(ctx, req.ID)                 // 第一刀: 写库前先删
err = NewProductDao(ctx).UpdateProduct(req.ID, product)
if err != nil {
    log.LogrusObj.Error(err); return
}
cache.DoubleDeleteAsync(req.ID, 0)                      // 第二刀: 500ms 后再删
emitProductChanged(ctx, req.ID, "update")               // outbox -> RMQ -> ES
\end{lstlisting}
\small \textbf{业务侧约定}：商家改价后\textbf{最长 500ms + Redis TTL}才能保证不出旧值；详情页 TTL=10min 时窗口实际由双删的 500ms 兜住。
\end{frame}

\begin{frame}{业务侧为什么接受"500ms 旧值"}
\begin{itemize}
  \item 强一致方案（2PC / 同步写）单次写要 50--200ms，且任何下游慢都会让"改价"按钮卡住。
  \item 客服话术："改价后 1 秒内可能仍显示老价格，刷新即可。"——这是和产品/运营事先谈好的口径，不是 bug。
  \item 反过来，如果业务允许 1s 误差，第二刀的 500ms 已经是过设计；改成 200ms 也行。
  \item 真正不能容忍旧价的场景是\textbf{下单页}：那里 \texttt{/orders/create} 必须从 DB 现读价格，不走 cache。
\end{itemize}
\keypoint{Cache Aside 是业务 SLA "最终一致 + $\le$ 1s 窗口" 的工程实现，不是教科书最优解。}
\srcnote{internal/product/service.go:271--292; repository/cache/product.go:64--74}
\end{frame}

% ============================================================
\section{HTTP cache：把流量挡在源站之外}
% ============================================================

\begin{frame}{为什么轮播图 24h 可以、价格 5 分钟都不行}
\textbf{核心问题}：每个接口 TTL 是\textbf{业务约定}决定的，不是技术决定的。
\vspace{1mm}
\begin{tabular}{@{}llll@{}}
\toprule
接口 & 业务变更频率 & 用户容忍延迟 & gomall TTL \\
\midrule
轮播图 \texttt{/carousels}      & 运营每天 1--2 次   & "下次刷首页可见即可" $\sim$ 10min  & 300s \\
类目 \texttt{/category/list}    & 几乎不变           & 几小时              & 300s \\
列表 \texttt{/product/list}     & 上架/下架几分钟一次 & "翻页内一致即可"   & 30s  \\
详情 \texttt{/product/show}     & 商家随时改价        & "刷新就能看到新价" & 60s  \\
下单页价格 (\texttt{/orders/create}) & 改价即生效         & \textbf{零}        & \textbf{不挂} \\
\bottomrule
\end{tabular}
\vspace{2mm}
\small 同样是商品域，下单页\textbf{绝对不能挂 cache}：用户付款瞬间必须用 DB 现读价格，否则发生"页面显示 ¥99 但实际扣 ¥199" = 大型客诉 + 资损。\\
轮播图 24h 也可以，但 gomall 取 5 分钟是\textbf{留出运营紧急下架的窗口}（违规图片、错图等）。
\srcnote{routes/router.go:39--45 公开 GET 挂 cache / 下单写接口 :81 走 authed 中间件链不挂 HTTPCache}
\end{frame}

\begin{frame}{四层缓存漏斗}
\begin{tikzpicture}[node distance=4mm and 0mm]
  \node[node-box,minimum width=70mm,node-ok] (b) {浏览器 cache (Cache-Control max-age)};
  \node[node-box,minimum width=58mm,node-ok,below=of b] (cdn) {CDN / 反向代理 (ETag + If-None-Match)};
  \node[node-box,minimum width=46mm,node-warn,below=of cdn] (rd) {Redis (product:detail:\{id\})};
  \node[node-box,minimum width=34mm,node-bad,below=of rd] (db) {MySQL (product PK)};
  \node[font=\tiny,right=2mm of b]   {\textgreater 60\% 流量在此截断};
  \node[font=\tiny,right=2mm of cdn] {再截 15--20\%};
  \node[font=\tiny,right=2mm of rd]  {再截 99\% 剩余};
  \node[font=\tiny,right=2mm of db]  {兜底回源};
\end{tikzpicture}
\vspace{2mm}
\small 业务侧的数字：上一档 cache 卸下的流量，比下一档便宜一个数量级。
浏览器 304 没有源站成本，CDN 一次回源覆盖一千个用户，Redis 一次回填覆盖 10min 内所有同商品请求。
\srcnote{routes/router.go:39--45 四档挂载点}
\end{frame}

\begin{frame}{每层命中率怎么算}
\begin{center}
\footnotesize
\begin{tabular}{@{}lp{4.4cm}p{5.0cm}@{}}
\toprule
层 & 命中率定义 & 度量口径 \\
\midrule
浏览器 & \texttt{304 / (200+304)}（按 referer 维度） & 前端埋点 \texttt{performance.getEntriesByType} \\
CDN    & \texttt{(edge\_hit + 304) / total\_req}    & CDN 厂商日志 \texttt{cache\_status=HIT} \\
Redis  & \texttt{1 - cache\_miss / total\_req}      & \texttt{INFO stats} 中 \texttt{keyspace\_hits / (hits+misses)} \\
DB     & 反向口径：\texttt{1 - 上三层总命中率}      & 由 prometheus \texttt{db\_qps / req\_qps} 反推 \\
\bottomrule
\end{tabular}
\end{center}
\vspace{2mm}

{\small 业务侧不会单独看 Redis 命中率，会看\textbf{端到端漏斗}：浏览器 60\%、CDN 在剩下 40\% 里再截 20\%（绝对值 8\%）、Redis 在剩下 32\% 里再截 99\%（绝对值 31.7\%），DB 最终承担 \textbf{0.32\%} 的原始流量。10k QPS 进来，DB 实际只看到 32 QPS。}
\srcnote{Prometheus 指标：\texttt{redis\_hits\_total} / \texttt{http\_requests\_total\{code="304"\}}}
\end{frame}

\begin{frame}{命中率链式公式与业务含义}
\keypoint{设各层命中率为 $h_1, h_2, h_3$，DB 流量 $= Q \cdot (1-h_1)(1-h_2)(1-h_3)$。}
\begin{itemize}
  \item \textbf{乘性效应}：任一层命中率掉 10\%，DB 流量翻倍不止。Redis 99\% $\to$ 90\% 时 DB 流量从 0.32\% 涨到 3.2\%，单机 MySQL 直接吃不消。
  \item \textbf{业务对齐}：浏览器层 60\% 不是技术指标，是\textbf{产品约定 max-age=60s} 换来的。产品坚持"改价 1s 内必须可见"时这层掉到 0。
  \item \textbf{容量规划反推}：先定 DB 能扛多少 QPS，倒推上游每层最低命中率，再倒推 TTL。
\end{itemize}
\srcnote{结合 stressTest/REPORT.md DB 单点 6 万级 QPS 与 product/show 62K RPS 的对齐}
\end{frame}

\begin{frame}{CDN 选型与命中率掉档的运维剧本}
\begin{tabular}{@{}llll@{}}
\toprule
方案 & 节点 / 成本 & 命中率 KPI & 适用 \\
\midrule
阿里云 CDN     & 国内 2800+ / 0.2 元/GB & 95\%+ & 国内电商主流 \\
Cloudflare    & 全球 300+ / 0.05 USD/GB & 92\%+ & 出海 / 多区域 \\
自建 ATS       & 0 边际，前期 1M+      & 看运维 & 月流量 \textgreater 10 PB \\
回源直发       & 0 + 源站带宽费        & N/A   & 试点 / 内网（gomall 当前）\\
\bottomrule
\end{tabular}
\vspace{2mm}
\small \textbf{掉档剧本}：浏览器突降 = 前端发版改了 URL query，回滚或加 \texttt{Vary} 白名单；CDN 突降 = 回源带 cookie，\texttt{/product/show} 必须无 session；Redis 突降 = \texttt{evicted\_keys} 在涨，扩容或调 \texttt{allkeys-lfu}；DB QPS 突涨但 Redis 没变 = 新接口漏挂 cache，grep \texttt{HTTPCache} 白名单。
\srcnote{Cloudflare cache analytics doc / 阿里云 CDN 计费规则 / Prometheus alert \texttt{cache\_hit\_ratio < 0.95 for 5m}}
\end{frame}

\begin{frame}{ETag 304 流程}
\begin{tikzpicture}[node distance=10mm]
  \foreach \name/\x in {Client/0, Gin/60, Handler/120} {
    \node[node-box,minimum width=16mm] at (\x mm,0) (\name) {\name};
    \draw[lifeline] (\x mm,-2mm) -- (\x mm,-58mm);
  }
  \node[font=\tiny\bfseries,align=left] at (-2mm,-8mm) {首次};
  \draw[arrow] (0,-12mm)  -- node[above,font=\tiny]{GET /product/show?id=42} (60mm,-12mm);
  \draw[arrow] (60mm,-18mm) -- node[above,font=\tiny]{ProductShow} (120mm,-18mm);
  \draw[arrow] (120mm,-24mm) -- node[above,font=\tiny]{200 + body} (60mm,-24mm);
  \draw[arrow] (60mm,-30mm) -- node[above,font=\tiny]{200 + ETag W/"a1b2..." + Cache-Control} (0,-30mm);
  \node[font=\tiny\bfseries,align=left] at (-2mm,-38mm) {再次};
  \draw[arrow] (0,-42mm)  -- node[above,font=\tiny]{GET ... + If-None-Match: W/"a1b2..."} (60mm,-42mm);
  \draw[arrow-async] (60mm,-48mm) -- node[above,font=\tiny]{(本地不算 ETag，直接 304)} (60mm,-48mm);
  \draw[arrow] (60mm,-54mm) -- node[above,font=\tiny]{304 (no body)} (0,-54mm);
\end{tikzpicture}
\srcnote{middleware/httpcache.go:53--99 ETag 比对与 304 分支}
\end{frame}

\begin{frame}[fragile]{ETag 中间件关键代码}
\begin{lstlisting}[language=Go,firstnumber=78,
  caption={\srcnote{middleware/httpcache.go:78--99 比对 + 304}}]
etag := weakETag(buf.body.Bytes())              // SHA-256 前 16 字节 -> W/"..."

if c.Request.Header.Get("If-None-Match") == etag {
    h := original.Header()
    h.Set("ETag", etag); h.Set("Cache-Control", cacheControl)
    h.Del("Content-Length"); h.Del("Content-Type")
    original.WriteHeader(http.StatusNotModified) // 304: 不写 body
    return
}
h := original.Header()
h.Set("ETag", etag); h.Set("Cache-Control", cacheControl)
h.Set("Content-Length", strconv.Itoa(buf.body.Len()))
original.WriteHeader(http.StatusOK)
_, _ = original.Write(buf.body.Bytes())          // 200: 正常回写
\end{lstlisting}
\end{frame}

\begin{frame}{ETag 中间件逐行讲解 + 为什么 POST 不挂}
\begin{itemize}
  \item \textbf{弱 ETag (\texttt{W/"..."})}：只承诺"语义等价"，不承诺字节级等价。压缩/分块都不影响命中，业务上更适合 JSON。
  \item \textbf{先 \texttt{c.Next()} 再比对}：handler 算完才能拿 body，中间件用 \texttt{cacheBuffer} 包 \texttt{ResponseWriter}，代价是多一份内存。
  \item \textbf{只对 \texttt{status=200} 生效}：4xx / 5xx 直接透传，业务码 70001 (限流) 不会被错误地长期缓存。
  \item \textbf{304 后删 \texttt{Content-Length / Content-Type}}：标准要求，否则上游 nginx / CDN 把空 body 算成 0 字节但仍发头，浏览器会拒收。
  \item \textbf{POST/PUT 不挂}：写接口挂 cache 等于"重复触发同一动作"，与幂等冲突；\texttt{POST /orders/create} 命中 cache 会让用户以为下单成功但实际没打到后端，对账少一笔。
\end{itemize}
\srcnote{middleware/httpcache.go:47--99 + 方法判断 :54--58}
\end{frame}

\begin{frame}{TTL 怎么定：业务节奏决定}
\begin{tabular}{@{}lrll@{}}
\toprule
接口 & max-age & 业务理由 & 风险 \\
\midrule
\texttt{/product/show}    & \textbf{60s}  & 单品价格/库存可能改  & 改价 1min 内显示老价 \\
\texttt{/product/list}    & 30s           & 列表翻页频繁，改卖家少 & 上架 30s 内不可见 \\
\texttt{/category/list}   & \textbf{300s} & 类目几乎不变         & 几乎无 \\
\texttt{/carousels}       & 300s          & 运营每天换一次轮播图  & 几乎无 \\
\bottomrule
\end{tabular}
\vspace{3mm}
\small 数字不是抄的，是和\textbf{产品同学}对齐过的"用户能容忍多久看到旧值"。\\
60s 是为了和 ProductDetailTTL (10min) 解耦：即使 Redis 还是旧值，HTTP 层 60s 后会强制重新协商一次。
\srcnote{routes/router.go:39 (30s) / :40 (60s) / :44 (300s) / :45 (300s)}
\end{frame}

% ============================================================
\section{三态故障：击穿 / 穿透 / 雪崩}
% ============================================================

\begin{frame}{三态故障的真实事故剧本}
\textbf{为什么这一节重要}：缓存三态故障是详情页能不能上 99.9\% 可用率的天花板。
\vspace{1mm}
\begin{description}
  \item[击穿事故] 双 11 0 点主推品 \texttt{product:detail:42} 因为运维误操作 \texttt{DEL}，瞬间 5000 个并发用户全打 DB，DB 连接池打满，整站详情页 503 持续 90s。损失：当 1 分钟 GMV $\sim$ 200 万。
  \item[穿透事故] 黑产脚本 100 QPS 扫 \texttt{?id=99999999..}，DB 每次 \texttt{record not found} 不回填 cache，DB CPU 100\%，影响下单页 \texttt{SELECT FOR UPDATE} 拿不到连接，订单 RT 涨到 5s。
  \item[雪崩事故] 商家批量上架 1000 SKU 后整 10 分钟集中过期，T+10min 那一秒 1000 个并发回源，DB QPS 瞬时尖峰 1000，触发主从延迟告警。
\end{description}
\vspace{2mm}
\small 三个事故对客户都是\textbf{"为什么平时秒开的页面突然 5 秒打不开"}，对 SRE 是\textbf{三种不同的处置 SOP}。区分清楚是基本功。
\srcnote{对应代码：internal/product/service.go:50--58 SETNX 单飞 (击穿) / 规划 bloom.go (穿透) / 规划 jitter (雪崩)}
\end{frame}

\begin{frame}{三种缓存故障的区别}
\begin{tabular}{@{}lll@{}}
\toprule
名称 & 触发条件 & 后果 \\
\midrule
\textbf{击穿} & 单个热 key 过期，瞬间 N 个请求回源 & 单点 DB 被打爆 \\
\textbf{穿透} & 请求查\textbf{不存在}的 key，cache 永远 miss & 每次都打 DB，CPU 100\% \\
\textbf{雪崩} & 大量 key 同一时刻过期 & 整片回源，DB 链接池打满 \\
\bottomrule
\end{tabular}
\vspace{3mm}
\small 三态共同点：cache miss 引起 DB 流量\textbf{阶跃}上涨。区别在\textbf{key 的范围}（单 key / 不存在的 key / 多 key）和\textbf{时间分布}（瞬时 / 持续 / 同步）。
对应防护手段三种：\textbf{SETNX 单飞} / \textbf{布隆过滤器} / \textbf{TTL 抖动}。下面分别讲。
\srcnote{对应 internal/product/service.go:50--58 SETNX + 后续准备引入的布隆过滤器和 TTL jitter}
\end{frame}

\begin{frame}{击穿：SETNX 单飞回源}
\begin{tikzpicture}[node distance=10mm]
  \foreach \name/\x in {Req1/0, Req2/30, Req3/60, Redis/100, DB/150} {
    \node[node-box,minimum width=16mm] at (\x mm,0) (\name) {\name};
    \draw[lifeline] (\x mm,-2mm) -- (\x mm,-50mm);
  }
  \draw[arrow] (0,-8mm)  -- node[above,font=\tiny]{GET (miss)} (100mm,-8mm);
  \draw[arrow] (30mm,-12mm) -- node[above,font=\tiny]{GET (miss)} (100mm,-12mm);
  \draw[arrow] (60mm,-16mm) -- node[above,font=\tiny]{GET (miss)} (100mm,-16mm);
  \draw[arrow] (0,-22mm)  -- node[above,font=\tiny]{SETNX lock OK} (100mm,-22mm);
  \draw[arrow-bad] (30mm,-26mm) -- node[above,font=\tiny\color{cBad}]{SETNX FAIL, sleep 50ms} (100mm,-26mm);
  \draw[arrow-bad] (60mm,-30mm) -- node[above,font=\tiny\color{cBad}]{SETNX FAIL, sleep 50ms} (100mm,-30mm);
  \draw[arrow] (0,-36mm)  -- node[above,font=\tiny]{回源} (150mm,-36mm);
  \draw[arrow] (150mm,-42mm) -- node[above,font=\tiny]{回填 cache + DEL lock} (100mm,-42mm);
  \draw[arrow-async] (30mm,-46mm)  -- node[above,font=\tiny]{retry GET (hit)} (100mm,-46mm);
\end{tikzpicture}
\srcnote{internal/product/service.go:50--58 + repository/cache/product.go:56--62}
\end{frame}

\begin{frame}{穿透：布隆过滤器挡住"不存在"的 id}
\begin{itemize}
  \item 攻击场景：脚本爬虫扫 \texttt{/product/show?id=99999999}，DB 返回 \texttt{record not found}，cache 不会回填 nil。下一次请求继续打 DB。
  \item 直觉修法 1：cache 回填 \texttt{nil} 标记 + TTL 60s。代价：脏数据需要等过期，正常 id 上架后 60s 不可见。
  \item 直觉修法 2：DB 前置布隆过滤器，启动时灌入所有合法 \texttt{product\_id}，新建商品时增量 \texttt{BF.ADD}。99.99\% 不存在的 id 在 Redis 层直接拒绝。
  \item gomall 当前没接，但已经预留接口：\texttt{cache.ProductExists(id)} 在 ProductShow 入口可加一行短路。
\end{itemize}
\srcnote{规划中：repository/cache/bloom.go (Redis 4.0+ \texttt{BF.RESERVE product:bf})}
\end{frame}

\begin{frame}[fragile]{布隆过滤器伪代码}
\begin{lstlisting}[language=Go,firstnumber=1,
  caption={\srcnote{规划：repository/cache/bloom.go BF.MEXISTS 用法}}]
// ProductExists 用布隆过滤器快速判 id 是否可能存在
// 误判率 0.01%; 误判方向只可能是 false-positive (说存在但实际不存在)
// 不会出现 false-negative (说不存在但实际存在), 业务侧零脏读
func ProductExists(ctx context.Context, id uint) (bool, error) {
    res, err := RedisClient.Do(ctx, "BF.EXISTS", "product:bf", id).Bool()
    if err != nil { return true, err }  // BF 故障时放行, 走原回源链路
    return res, nil
}

// ProductShow 入口处加短路
exists, _ := cache.ProductExists(ctx, req.ID)
if !exists {
    return nil, e.ErrProductNotFound   // 0.01ms 直接拒绝
}
\end{lstlisting}
\end{frame}

\begin{frame}{雪崩：TTL 抖动与多级兜底}
\begin{itemize}
  \item 触发场景：商家批量上架 1000 个 SKU，全部用默认 \texttt{ProductDetailTTL = 10min}。10 分钟后\textbf{同一秒}集中过期，回源瞬时打 1000 次 DB。
  \item 修法 1：TTL \textbf{加随机抖动}。写入时 \texttt{ttl = base + rand(0, 0.2*base)}，把过期时间打散到 2 分钟窗口。
  \item 修法 2：\textbf{多级 TTL}。L1 cache 10min、L2 cache 1h，L1 过期回 L2，L2 过期回 DB。gomall 暂未启用 L2，路线图上是 Redis cluster 上的 \texttt{ro-replica}。
  \item 修法 3：\textbf{熔断兜底}。DB 拒连时整链路降级返回上一版 cache（哪怕过期），客服话术接受"价格延迟 1 分钟"。
\end{itemize}
\srcnote{规划：repository/cache/product.go SetProductDetail 加 jitter 实现}
\end{frame}

\begin{frame}{第二刀间隔的取舍}
\begin{tabular}{@{}lll@{}}
\toprule
第二刀间隔 & 覆盖率 & 副作用 \\
\midrule
100ms & 覆盖 $\sim$ 80\% 回填窗口 & 改价后体感最快，但漏窗仍有 \\
\textbf{500ms} & 覆盖 $\sim$ 99\% 回填窗口 & gomall 当前值，平衡点 \\
1s    & 覆盖 $\sim$ 99.9\% 回填窗口 & 改价后用户最长等 1s + TTL \\
3s    & 几乎全覆盖 & 用户体感差，且改价高峰时 goroutine 堆积 \\
\bottomrule
\end{tabular}
\vspace{3mm}
\small 数字怎么来：取决于\textbf{读路径回填耗时}的 P99。gomall 单次回填 = DB 查询 (3ms) + JSON marshal (1ms) + Redis SET (1ms) $\approx$ 5ms。考虑 GC / 抖动 / 网络重传，P99 不超过 100ms；500ms 留 5 倍余量。
\srcnote{repository/cache/product.go:16 ProductDelayInterval = 500ms}
\end{frame}

\begin{frame}{热 key：探测 + 三种拆分策略}
\small
\textbf{现象}：单分片 CPU 100\%、出带宽打满，其他分片闲；典型是秒杀 / 双 11 主推品 \texttt{product:detail:\{id\}}。
\vspace{1mm}
\textbf{探测三法}：(1) \texttt{redis-cli --hotkeys} LFU 采样、暂停 4s，离线用；(2) \texttt{slowlog + MONITOR} 抽样 1\% 按 key 聚合；(3) Gin 层 \texttt{sync.Map} 埋点 dump Prometheus、超 1000 QPS 告警。gomall 当前用方法 3。
\vspace{1mm}
\textbf{拆分策略（按读写比选）}：
\begin{description}
  \item[本地 process cache] Go 进程内 \texttt{LRU(100,5s)} 绕过 Redis；一致窗口 5s，适合\textbf{读写 \textgreater 1000:1} 的商品详情。
  \item[Key 分片] \texttt{product:detail:42:shard0..N} 读随机写全更，N=10 单分片 QPS 降 10 倍，写放大。
  \item[读写分离] 热 key 走 \texttt{ro-replica}，一致窗口 = 主从同步 $\sim$ 10ms；秒杀库存读写 1:1 必须分片。
\end{description}
\srcnote{规划：middleware/hotkey.go + Prometheus \texttt{product\_view\_total} / Twitter Manhattan / Facebook MCRouter}
\end{frame}

% ============================================================
\section{反例：product/list 全表 COUNT}
% ============================================================

\begin{frame}{业务侧视角：为什么 list 是真正的瓶颈}
\begin{itemize}
  \item C 端用户：列表页是\textbf{发现商品}的主入口，比详情更上游；list 慢 = 用户根本进不来详情。
  \item 商家：上架新 SKU 后\textbf{30s 内可见}（HTTP TTL），但分页到第 5 页用户根本看不到 = 商家流量分配不公平。
  \item 运营：类目页 "全部 1{,}245 件" 的 total 是\textbf{产品需求} —— 但用户从来不在乎是 1245 还是约 1200，是过设计。
  \item SRE：\texttt{product/list} 24.5 RPS / p95 2.5s 是\textbf{1M 行表上 COUNT(*) 不可避免的代价}，不是 bug。要么改产品需求，要么换实现。
\end{itemize}
\vspace{1mm}
\textbf{业务取舍排序}：
\begin{enumerate}
  \item 干掉 total，前端改无限滚动（最优，但要拉前端 PM 一起对齐）
  \item 显示"约 1{,}000 件"用近似计数（中等，前端文案改一处）
  \item 物化视图 + outbox 异步刷（最重，QPS \textgreater 1k 才划算）
\end{enumerate}
\srcnote{stressTest/REPORT.md \S 3 商品列表反模式 + internal/product/repo.go:58--64}
\end{frame}

\begin{frame}{反例数字：24.5 RPS / p95 2.5s}
\begin{tabular}{@{}lrr@{}}
\toprule
接口 & 实测 RPS & p95 \\
\midrule
\texttt{/ping} 基线          & 64{,}254 & 3.51ms \\
\texttt{/product/show}      & 62{,}226 & 3.01ms \\
\texttt{/product/list} (反例) & \textbf{24.5}    & \textbf{2.50s} \\
\bottomrule
\end{tabular}
\vspace{2mm}
\small 同样是商品域的 GET 接口，吞吐差 \textbf{2{,}500 倍}，p95 差 \textbf{800 倍}。\\
区别只有一个：\texttt{ShowProduct} 走主键，\texttt{ListProducts} 多了一个 \texttt{SELECT COUNT(*)}。
\srcnote{stressTest/REPORT.md \S 各链路压测结果（行 product/list）}
\end{frame}

\begin{frame}{扫表示意：1M 行的 COUNT 怎么炸的}
\begin{tikzpicture}
  \draw[draw=cMain,thick] (0,0) rectangle (10,0.8);
  \foreach \i in {0,1,2,3,4,5,6,7,8,9} {
    \draw[draw=cGray,dashed] (\i,0) -- (\i,0.8);
  }
  \node[font=\tiny,below] at (5,-0.1) {product 表 956K 行 / 364 MB};
  \node[font=\tiny,above] at (5,0.9) {\texttt{SELECT COUNT(*) FROM product WHERE category\_id=?}};
  \fill[cBad!40] (0,0) rectangle (10,0.8);
  \node[font=\small\bfseries,white] at (5,0.4) {全表扫: 每次请求 $\sim$ 2 s};
  \node[font=\tiny\color{cBad},right] at (10.1,0.4) {50 VU * 2s $\to$ 24.5 RPS};
\end{tikzpicture}
\vspace{6mm}
\begin{itemize}
  \item \texttt{Count(\&total)} 在没有 \texttt{category\_id} 索引时退化为全表扫。
  \item 即使有索引，InnoDB 的 \texttt{COUNT(*)} 不能像 MyISAM 那样常量返回，仍要走索引页。
  \item 业务侧问的从来不是"商品一共多少件"，而是"我这一页有多少件"——total 大多是\textbf{产品需求过度}。
\end{itemize}
\srcnote{internal/product/repo.go:58--64 CountProductByCondition}
\end{frame}

\begin{frame}[fragile]{反例代码}
\begin{lstlisting}[language=Go,firstnumber=58,
  caption={\srcnote{internal/product/repo.go:58--64}}]
// CountProductByCondition 根据情况获取商品的数量
func (d *ProductDao) CountProductByCondition(condition map[string]interface{}) (total int64, err error) {
    err = d.DB.Model(&Product{}).
        Where(condition).Count(&total).Error    // 1M 行表 + 无 category_id 索引 -> 全表扫
    return
}
\end{lstlisting}
\vspace{2mm}
\textbf{service 调用点}：\srcnote{internal/product/service.go:195--204 ListProducts 同时调 ListProductByCondition + CountProductByCondition}
\end{frame}

\begin{frame}{三种修法对业务的代价}
\begin{tabular}{@{}lll@{}}
\toprule
方案 & 代价 & 业务可见变化 \\
\midrule
建 \texttt{category\_id} 索引 & 写入慢 5--10\% & 几乎无 \\
近似计数 (\texttt{TABLE STATUS}) & total 有 5\% 误差 & 显示 "约 1{,}000 件" \\
物化视图 (\texttt{category\_count}) & 异步刷新延迟 & 数字偶有 5min 滞后 \\
干掉 total，前端无限滚动 & 产品改交互 & 翻页器消失 \\
\bottomrule
\end{tabular}
\vspace{3mm}
\small 后端常常想直接选"建索引"，但和产品/运营对齐后会发现\textbf{无限滚动}才是真正合理的选择：用户分页到第 50 页的概率 $<$ 0.1\%，业务上没人在乎 total。
\srcnote{stressTest/REPORT.md \S 3 商品列表反模式 改进方向}
\end{frame}

\begin{frame}{物化视图 / 无限滚动：选哪个}
\begin{tabular}{@{}lllr@{}}
\toprule
方案 & 一致性 & 复杂度 & 查询耗时 \\
\midrule
\texttt{(category, on\_sale)} 复合索引 & 强 & 加索引一行 & 50--200ms \\
物化视图 \texttt{product\_count\_by\_cat} & 5min 最终一致 & 中（需 outbox） & \textless 1ms \\
近似计数 \texttt{TABLE\_ROWS} & 弱（10\% 误差） & 极低 & \textless 1ms \\
无限滚动（cursor） & 不需要 total & 中（前端改） & \textless 1ms \\
\bottomrule
\end{tabular}
\vspace{2mm}
\small 物化视图临界点：\textbf{QPS \textgreater 1k \&\& 类目稳定}。gomall 当前 24.5 RPS 没到，先建索引。\\
无限滚动是最优解：用户翻第 50 页概率 \textless 0.1\%，total 价值极低。\texttt{order/list} 已用游标 + 缓存把 QPS 从 8.3 拉到 \textbf{58{,}406}（\textasciitilde 7000 倍），商品列表可复用同样模板：cursor = \texttt{(sort\_key, id)} 二元组。
\srcnote{游标分页模板已在 stressTest/REPORT.md \S 2 验证}
\end{frame}

\begin{frame}{容量规划：1k / 10k / 100k QPS 链路资源估算}
\begin{tabular}{@{}lrrr@{}}
\toprule
入口 QPS & 1{,}000 & 10{,}000 & 100{,}000 \\
\midrule
浏览器卸载（60\%） & 400 到源 & 4{,}000 到源 & 40{,}000 到源 \\
CDN 卸载（20\% 再截） & 320 到源站 & 3{,}200 到源站 & 32{,}000 到源站 \\
源站 Redis（99\% 命中） & 3.2 到 DB & 32 到 DB & 320 到 DB \\
DB 实际负载 & \textless 1\% 单机 & 10\% 单机 & 接近上限 \\
源站 CPU（按 50K RPS/核） & 0.01 核 & 0.1 核 & 1 核 \\
\bottomrule
\end{tabular}
\vspace{2mm}
\small 100k QPS 才把 DB 推到 320 QPS，远低于单机 6 万级 PK 上限。\textbf{单机够用}，所以 gomall 在百万级 DAU 之前不需要分库分表，只需要把 cache 命中率守住 99\%。
\srcnote{基线数字来自 stressTest/REPORT.md /ping 64K RPS / product/show 62K RPS}
\end{frame}

% ============================================================
\section{业务码、客诉与监控口径}
% ============================================================

\begin{frame}{客服话术手册：商品展示侧 4 类典型客诉}
\begin{center}
\footnotesize
\begin{tabular}{@{}p{32mm}p{20mm}p{48mm}@{}}
\toprule
客诉描述 & 责任方 & 客服回复 + 升级路径 \\
\midrule
"价格不对，详情写 99 列表写 88" & 后端缓存层 & "刷新即可。1 分钟内仍不一致请截图反馈。"——开 P3 工单查 outbox / 双删 \\
"商品已经下架了首页还在推" & 运营 + 后端 & "我们立刻同步下架。"——运营 console DEL cache + 强推 outbox \\
"刚下单的商品库存数和实际不一致" & 见 deck 07 & "您的订单已锁库存，不会超卖。"——库存事件中心 \\
"详情页一直转圈打不开"     & SRE       & "活动太火爆，1 分钟后重试。"——查 70002 熔断 / DB QPS \\
\bottomrule
\end{tabular}
\end{center}
\vspace{2mm}

{\small 客服话术不是事后想的，是\textbf{写代码时就和客服 / 运营对齐的合同}。70001 / 70002 / 60002 三个状态码各自对应一句固定话术，让客服不用判断技术细节也能回应。}
\srcnote{pkg/e/code.go:53--58 业务码 + README \S 业务全景 客服角色行}
\end{frame}

\begin{frame}{各角色仪表盘：每人看什么}
\footnotesize
\begin{tabular}{@{}p{14mm}p{55mm}p{28mm}@{}}
\toprule
角色 & 关键指标 & 报警阈值 \\
\midrule
C 端 & 详情 p99 / 列表 p99 / 首屏 LCP & p99 \textgreater 200ms 5min \\
商家 & 上架 $\to$ 可见耗时 / 改价 $\to$ 可见耗时 & \textgreater 60s / \textgreater 1s \\
运营 & 轮播 CTR / 类目导流 GMV / 当日上架量 & CTR 同比 -20\% \\
客服 & 70001 / 70002 / 60002 计数 / 工单量 & 工单 \textgreater 10/min \\
SRE  & Redis 命中率 / DB QPS / 304 占比 / outbox lag & 命中率 \textless 95\% 5min \\
\bottomrule
\end{tabular}
\vspace{2mm}

\textbf{要点}：同一个 cache 命中率，SRE 看的是"我系统健不健康"，运营看的是"我商家上架速度被压了没"，是\textbf{同一份原始数据的不同业务视角}。Prometheus 一套指标支撑多个仪表盘。
\srcnote{Prometheus 指标 redis\_hits\_total / http\_requests\_total\{code="304"\} / outbox\_pending}
\end{frame}

\begin{frame}{业务码：60002 的真实含义}
\footnotesize
\begin{tabular}{@{}lll@{}}
\toprule
状态码 & 含义       & 客服话术 \\
\midrule
60001  & 幂等 token 无效 & "请刷新页面后重试" \\
60002  & 幂等正在处理   & "您的请求正在处理，1 秒后自动刷新" \\
70001  & 限流命中       & "活动太火爆，请稍后重试" \\
70002  & 熔断开启       & "支付服务暂忙，1 分钟后" \\
\bottomrule
\end{tabular}
\vspace{2mm}

注意 60002 不是\textbf{cache 命中}（HTTP 304 才是），而是\textbf{幂等命中}——同一 Idempotency-Key 在第一笔还没结束就再来一次。
两者对客服来说都是"无需重复操作"，话术接近；但监控口径不同：304 不写日志、不计错误率；60002 要计入"幂等保护次数"。
\srcnote{pkg/e/code.go:53--58 业务码定义}
\end{frame}

\begin{frame}{客诉拆解 + 业务侧最终目标}
\textbf{"为什么显示老价格"三种场景}：
\begin{description}
  \item[A] 改价后 $\le$ 500ms：双删第二刀还没到，让用户再刷一次。
  \item[B] 改价后 1--60s：浏览器/CDN 强 cache，教硬刷新 (\texttt{Cmd+Shift+R}) 或运营"强制下架重新上架"。
  \item[C] 改价后 $>$ 10min 仍旧值：双删失败或 outbox 没投递，开 P2 工单查 \texttt{outbox} 表。
\end{description}
\vspace{2mm}
\textbf{业务侧最终目标}：首屏 1s（浏览器 304 + CDN）/ 命中率 \textgreater 99\%（Redis + SETNX）/ 改价 1s 内可见（双删 + outbox）/ 翻页不炸（cursor + 无限滚动）。
\keypoint{商品展示不是"把数据从 DB 取出来"，是\textbf{把流量挡在 DB 之外}。}
\srcnote{internal/product/service.go:282--289 双删 + emitProductChanged 三联}
\end{frame}

\begin{frame}{商品展示路线图：从 P1 信息类到 P0 流量入口}
\begin{tabular}{@{}p{20mm}p{42mm}p{36mm}@{}}
\toprule
阶段 & 已做 & 路线图 \\
\midrule
M1 当前 & Cache Aside / SETNX / 双删 / HTTPCache / outbox 联动 ES & — \\
M2 下个迭代 & — & 布隆过滤器 \texttt{cache.ProductExists} / TTL jitter / 热 key 探测 \\
M3 商家增长 & BossID 字段就位 & 商家自助上架 console + 改价审计日志 \\
M4 大促准备 & 静默降级哲学 & 多级 cache (L1 process + L2 Redis cluster) / CDN 选型 \\
M5 个性化 & — & 详情底部"猜你喜欢"召回 + AB 框架 \\
\bottomrule
\end{tabular}
\vspace{2mm}
\small 业务等级演进：M1 详情属于 P1 信息类（99.9\% / p99 200ms），M4 之后双 11 主推品时段升级到\textbf{准 P0}（与下单同等级），届时要给 cache 加单分片冗余 + 客户端旁路缓存。
\srcnote{README.md \S 业务承诺(SLO) + \S 业务边界 + docs/architecture/feature-matrix.md}
\end{frame}

% ============================================================
\appendix
% ============================================================

\begin{frame}{面试 Q\&A（一）：业务取舍篇}
\small
\textbf{Q1}：为什么轮播图 24h cache 可以，价格 5 分钟都不行？\\
A：业务变更频率不同。运营改轮播是手动行为、用户对"今天看到昨天的图"零容忍度；价格是商家随时改的、用户对老价高度敏感（漏发优惠 / 资损客诉）。TTL 跟着\textbf{业务约定}走，不是技术决定。

\vspace{1mm}
\textbf{Q2}：商家改价 1s 内可见这条 SLA 是谁定的？\\
A：和产品 / 运营当面对齐过的合同。同步写 50--200ms 体验差且任何下游慢就卡按钮；500ms 双删 + TTL 留 5 倍余量。客服话术写明"1 秒内可能旧价，刷新即可"——出问题不算事故。

\vspace{1mm}
\textbf{Q3}：商家改价后用户客诉"显示老价格"，怎么定位 + 客服怎么说？\\
A：按时间窗口拆。$\le$ 1s 是双删窗口（正常，让用户刷新）；1--60s 是 HTTP cache（教硬刷新或运营强制下架重新上架）；$>$ 10min 是双删失败 / outbox 没投（开 P2 工单查 \texttt{outbox} 表）。

\vspace{1mm}
\textbf{Q4}：商家上架 1 万 SKU 后 DB 顶不住的事故剧本怎么避免？\\
A：核心是\textbf{缓存击穿防护}。运营把头部 SKU 引流 = 这些 \texttt{product:detail:\{id\}} 第一次访问全是 cache miss。SETNX 单飞保证同商品最多 1 个请求打 DB，剩余 sleep 50ms 后命中回填。没有 SETNX = 一万次并发 SELECT 同时打 DB = 整站详情页 503。
\end{frame}

\begin{frame}{面试 Q\&A（二）：角色与边界篇}
\small
\textbf{Q5}：商品展示这一节解决的是哪个角色的什么痛点？\\
A：C 端是转化漏斗最宽的一跳（详情页 RT +100ms = 跳出率 +2\%）；商家是\textbf{流量分配}（上架到可见 SLA = 60s）；运营是首页 GMV 杠杆（轮播加载慢 1s = 首页 GMV $-$2\%）；客服是话术口径；SRE 是 99\% 命中率守住 = DB 单机够用。

\vspace{1mm}
\textbf{Q6}：详情页这一节\textbf{不解决}什么？\\
A：不做 SKU 多规格、不做视频 / 直播、没 AB 框架、没个性化推荐、商品搜索召回在 deck 03、评价 / 售后在 deck 09 路线图阶段。摊清边界比堆 feature 值钱。

\vspace{1mm}
\textbf{Q7}：商家秒级改价、用户必须秒级看到，怎么办？\\
A：改\textbf{业务约定}，不改技术。要么 max-age=0 取消强缓存，要么走下单页直读 DB 兜底。"既要强缓存又要强一致"在分布式里没有解 —— 这是要去找产品谈的，不是工程师能拍板的。

\vspace{1mm}
\textbf{Q8}：缓存击穿、穿透、雪崩三个名词怎么区分？业务侧后果分别是？\\
A：击穿=单 key 过期瞬时并发回源 (SETNX 单飞) —— 双 11 主推品 \texttt{DEL} 误操作，1 分钟丢 200 万 GMV；穿透=查不存在的 key 一直 miss (布隆过滤器) —— 黑产扫 id 打 DB CPU 100\%；雪崩=多 key 同时过期 (TTL 抖动) —— 商家批量上架后 T+10min 集中过期。
\end{frame}

\begin{frame}{面试 Q\&A（三）：SLO 与监控篇}
\small
\textbf{Q9}：商品详情背的 SLO 是什么？什么时候算事故？\\
A：P1 信息类 99.9\% / p99 \textless 200ms，实测 p95 3ms 有 60 倍余量。事故触发：p99 \textgreater 200ms 持续 5min，或可用率跌破 99.9\%（年宕 8.76h 预算）。商品列表 24.5 RPS / p95 2.5s 当前 \textbf{不算事故}（P3 上限 1s 但 QPS 远低于阈值），属于"应该排期修"。

\vspace{1mm}
\textbf{Q10}：DB 实际承担多少 QPS 怎么算？哪层掉了会怎样？\\
A：链式公式 $Q \cdot (1-h_1)(1-h_2)(1-h_3)$。浏览器 60\% / CDN 20\% / Redis 99\%，10k QPS 进 DB 看到 32 QPS。\textbf{Redis 雪崩时}浏览器/CDN 没异常，总命中率仍 80\%+ 看着像没事，但 DB QPS 翻 10 倍。必须分层报警（\texttt{redis\_hit \textless 95\% for 5m}）。

\vspace{1mm}
\textbf{Q11}：\texttt{COUNT(*)} 四种修法业务上怎么选？谁能拍板？\\
A：QPS \textless 1k 加索引（后端独立决定）；要"约 1000 件"文案用近似计数（要拉产品对文案）；类目稳定且 QPS \textgreater 1k 上物化视图（架构师拍板）；无限滚动 = 最优解但要前端 PM + 产品三方对齐。\textbf{业务取舍优先级 \textgreater 技术选型}。

\vspace{1mm}
\textbf{Q12}：百万级 DAU 时需要分库分表吗？\\
A：不需要。100k QPS 入口、99\% 命中率，DB 实际看 320 QPS，远低于单机 6 万级 PK 上限（\texttt{stressTest/REPORT.md} 实测）。先把 cache 命中率守住，再考虑分片。\textbf{别为不存在的瓶颈做不存在的优化}。
\end{frame}

\begin{frame}{代码位置一览}
\scriptsize
\begin{description}
  \item[读路径 Cache Aside]\srcnote{internal/product/service.go:42--68 ProductShow}
  \item[写路径 双删]\srcnote{internal/product/service.go:271--292 ProductUpdate}
  \item[Cache 包装]\srcnote{repository/cache/product.go:30--74 Get/Set/Del/DoubleDeleteAsync}
  \item[SETNX 回源锁]\srcnote{repository/cache/product.go:55--62 TryProductLock}
  \item[Double Delete 异步]\srcnote{repository/cache/product.go:64--74 DoubleDeleteAsync 500ms}
  \item[HTTP cache 中间件]\srcnote{middleware/httpcache.go:47--106 HTTPCache + weakETag}
  \item[ETag 304 比对]\srcnote{middleware/httpcache.go:78--99 If-None-Match 分支}
  \item[路由挂载]\srcnote{routes/router.go:39--40 (show 60s / list 30s) + :44--45 (300s)}
  \item[反例：全表 COUNT]\srcnote{internal/product/repo.go:58--64 CountProductByCondition}
  \item[规划：布隆过滤器]\srcnote{规划 repository/cache/bloom.go BF.EXISTS}
  \item[规划：TTL jitter / 热 key 探测]\srcnote{规划 SetProductDetail + rand jitter / middleware/hotkey.go}
  \item[业务码]\srcnote{pkg/e/code.go:53--58 (60001/60002/70001/70002)}
  \item[Outbox 联动 ES]\srcnote{internal/product/service.go:257--265 emitProductChanged + :155/253/289 三处调用点}
  \item[业务全景 / SLO]\srcnote{README.md \S 业务全景 / 用户旅程 / 业务承诺(SLO) / 业务边界}
  \item[压测数字]\srcnote{stressTest/REPORT.md 链路压测表 + \S 3 反模式 + 基线 64{,}254 / 62{,}226 / 24.5 RPS}
\end{description}
\end{frame}

\end{document}
